\begin{theorem}[Norm filtration below and at the break]
\label{mod:ramification-normbelowbreak}
Let $K/F$ be a ramified cyclic extension of prime degree
$\ell=[K:F]$, let $t\in\mathbf N$ be its unique lower break, and assume
$f(K/F)=1$.  Fix an integral uniformizer $\pi_K$ together with the
explicit generator hypothesis
\[
 \operatorname{Algebra.adjoin}_{\mathcal O_F}\{\pi_K\}=\mathcal O_K.
\]
The residue-degree assumption is retained explicitly throughout; in
particular, it gives $e(K/F)=\ell$ and identifies the cardinalities of the
two residue fields and of their corresponding unit-graded pieces.

Write
\[
 \operatorname{gr}^iU_E=U_E^i/U_E^{i+1}.
\]
For two certified depths $a,b$, the predicate
\texttt{NormMapsUnitFiltration} records the containment
$N_{K/F}(U_K^a)\subseteq U_F^b$.  The norm is then restricted to these
subgroups and descended to quotient groups.  Numerator containment and
denominator containment are separate hypotheses, so representative
independence is part of the construction.  In particular, the exact graded
norm is the homomorphism
\[
 \overline N_i:\operatorname{gr}^iU_K\longrightarrow
                 \operatorname{gr}^iU_F,
 \qquad [u]\longmapsto[N_{K/F}(u)],
\]
defined only after the containments at both $i$ and $i+1$ have been
supplied.  Its representative formula is an equality of quotient classes,
not a choice of a representative.

The trace input needed to prove these containments is established here,
independently of the final norm-filtration theorem.  Put
\[
 D=(\ell-1)(t+1).
\]
Trace duality for the chosen integral generator proves, for every
$q\in\mathbf Z$,
\[
 \operatorname{Tr}_{K/F}(\mathfrak p_K^q)
 \subseteq
 \mathfrak p_F^{\lfloor(q+D)/\ell\rfloor},
\]
and proves surjectivity onto the displayed target lattice.  Together these
are the exact trace-ideal equality used in the norm corrections; neither
direction is imported from Theorem~\ref{thm:norm-filtration}.

The ramification homomorphism is constructed before its relation to the
norm is considered.  For $\sigma\in G_i$, the unit
\[
 \rho_{i,\pi_K}(\sigma)=\frac{\sigma(\pi_K)}{\pi_K}
\]
belongs to $U_K^i$, and changing $\sigma$ by an element of $G_{i+1}$
does not change its class in $\operatorname{gr}^iU_K$.  This gives an
injective homomorphism
\[
 \overline\rho_i:G_i/G_{i+1}\hookrightarrow
                  \operatorname{gr}^iU_K.
\]
At $i=0$ this is the multiplicative residue coordinate
$\overline{\sigma(\pi_K)/\pi_K}$.  At positive depth, under the usual
additive identification of the graded line, the same quotient class is
the coordinate
\[
 \overline{\frac{\sigma(\pi_K)-\pi_K}{\pi_K^{i+1}}}.
\]
Its kernel is exactly $G_{i+1}$, and
$N_{K/F}(\rho_{i,\pi_K}(\sigma))=1$ exactly.  At a positive break, the
faithful positive-depth coordinate also proves
$\operatorname{char}k_F=\ell$ rather than taking wildness as an additional
hypothesis.

For $t>0$, the exact elementary-symmetric expansion supplies the
representative formulas needed for the graded maps.  If
$v_K(x)\geq i<t$, then
\[
 N_{K/F}(1+x)-1-N_{K/F}(x)\in\mathfrak p_F^{i+1};
\]
thus, for $i>0$, the additive coordinate of $\overline N_i$ is represented
by the terminal norm term; at $i=0<t$ the same calculation is the residue
Frobenius.  At the positive critical depth,
\[
 N_{K/F}(1+x)-1-\operatorname{Tr}_{K/F}(x)-N_{K/F}(x)
 \in\mathfrak p_F^{t+1}
 \qquad(v_K(x)\geq t),
\]
so precisely the trace and terminal norm terms survive.  The separate
deeper estimate controls all nontrace coefficients for successive lifting.
In the tame branch $t=0$, norm preserves $U^0$ and $U^1$; the critical map
is the honest residue-unit norm, and a corresponding trace congruence at
depths $\ell(n+1)$ drives the lifting above $U^1$.

Strictly below the break, for every $0\leq i<t$, the exact graded norm has
\[
 \ker(\overline N_i)=1,
 \qquad
 \operatorname{im}(\overline N_i)=\operatorname{gr}^iU_F,
\]
and the pair
$G_i/G_{i+1}\longrightarrow\operatorname{gr}^iU_K
\longrightarrow\operatorname{gr}^iU_F$ is exact: below the break
$G_i=G_{i+1}=G$, so its first term and the graded norm kernel are both
trivial.
The graded norm hence induces a multiplicative equivalence
\[
 \operatorname{gr}^iU_K\simeq\operatorname{gr}^iU_F.
\]
Injectivity follows from the representative congruence; surjectivity then
uses the equality of the finite source and target cardinalities.  Splicing
exactly the indices
\[
 i=r,r+1,\ldots,t-1
\]
gives
\[
 U_K^r/U_K^t\simeq U_F^r/U_F^t
 \qquad(0\leq r\leq t).
\]
The endpoint $i=t$ is deliberately excluded from this splice.  The same
finite induction produces both a representative modulo $U_F^t$ and the
reverse membership test: a norm lying in $U_F^t$ forces a source unit,
initially in $U_K^r$, to lie in $U_K^t$.  After matching valuations using
$f(K/F)=1$, these two statements give the corresponding norm equivalence
$K^\times/U_K^t\simeq F^\times/U_F^t$.

At the critical depth, numerator and denominator preservation are proved
at the separate endpoints $t$ and $t+1$, and the resulting map is
\[
 \overline N_t:U_K^t/U_K^{t+1}\longrightarrow
                 U_F^t/U_F^{t+1}.
\]
For both $t=0$ and $t>0$, the exact sequence at the middle term is
\[
 G_t/G_{t+1}\xrightarrow{\ \overline\rho_t\ }
 \operatorname{gr}^tU_K\xrightarrow{\ \overline N_t\ }
 \operatorname{gr}^tU_F,
 \qquad
 \operatorname{im}(\overline\rho_t)=\ker(\overline N_t),
\]
with $\overline\rho_t$ injective.  Hence
\[
 \#\ker(\overline N_t)=\ell,
 \qquad
 \#\bigl(\operatorname{gr}^tU_F/
                 \operatorname{im}(\overline N_t)\bigr)=\ell.
\]
For $t>0$, equivalently,
\[
 \#\operatorname{im}(\overline N_t)\,\ell=\#k_F.
\]
The cokernel is the displayed quotient by the actual norm range, so its
direction is fixed explicitly; its cardinality follows from the proved
exact kernel and equality of the finite graded cardinalities, not from an
assumed graded-kernel statement.

For an arbitrary positive lower break $t$---without presenting it as an
endpoint $s+1$---put $\pi_F=N_{K/F}(\pi_K)$ and use the normalized additive
coordinates on the two critical graded lines.  If $g$ is the chosen generator
of the cyclic Galois group, define
\[
 c=\overline{\frac{\operatorname{Tr}_{K/F}(\pi_K^t)}{\pi_F^t}}
 \quad\text{and}\quad
 \lambda=(k_F\simeq k_K)^{-1}
   \left(\overline{\frac{g(\pi_K)-\pi_K}{\pi_K^{t+1}}}\right).
\]
The critical norm polynomial in these coordinates is then proved exactly:
\[
 P_t(z)=z^\ell+c z
       =z^\ell-\lambda^{\ell-1}z,
 \qquad c=-\lambda^{\ell-1}.
\]
Thus the polynomial, its linear coefficient, and its ramification parameter
are reusable at every positive break, rather than being tied to either of the
later endpoint cases.  The same construction fixes the cyclic enumeration
$\mathbf Z/\ell\mathbf Z\simeq\operatorname{Gal}(K/F)$ by sending $1$ to
$g$ and proves that the actual residue displacement of the class $j$ is
$j\lambda$ in $k_K$.

The reverse critical-kernel inclusion is proved by an explicit
Hilbert--90 correction.  A representative whose norm is only trivial
modulo $U_F^{t+1}$ is first divided by an element of $U_K^{t+1}$ having
the same exact norm.  This yields a representative $w\in U_K^t$ with
$N_{K/F}(w)=1$ and with the original class in
$U_K^t/U_K^{t+1}$.  Only after this exact correction is Hilbert 90 applied,
in the fixed orientation
\[
 \frac{y}{g(y)}=w,
\]
where $g$ is the chosen generator of the cyclic Galois group.  Removing the
uniformizer power from $y$ shows that the class of $w$ is a power of the
ramification coordinate.  The tame and positive-break corrections are
proved separately and then combined in the principal kernel theorem.

Finally, the lifting depths used in that correction are explicit.  For a
positive break, a target error in $U_F^{t+1+n}$ is corrected one level
deeper by a source element at depth
\[
 t+\ell n+(\ell-1)+1=t+\ell(n+1)
   =\Psi_{t,\ell}(t+1+n),
\]
leaving an error in $U_F^{t+1+n+1}$.  In the tame case the corresponding
source depth is
$\ell n+(\ell-1)+1=\ell(n+1)=\Psi_{0,\ell}(1+n)$.
The monotonicity and divergence of these depths, completeness, and
continuity of norm invoke the lifting mechanism of
module~\ref{mod:localfield-successivelifting}; they give exact norm
representatives for $U_F^{t+1}$ in $U_K^{t+1}$ (and for $U_F^1$ in
$U_K^1$ when $t=0$).  Notice the endpoint distinction
\[
 \Psi_{t,\ell}(t)=t,
 \qquad
 \Psi_{t,\ell}(t+1)=t+\ell,
\]
and the different roles of $t+1$ and $t+\ell$: the critical graded
denominator is $U^{t+1}$, whereas the first positive-break correction is
chosen at the Herbrand depth $t+\ell$.

\LeanFile{LanglandsFirstMainLemma/Ramification/NormBelowBreak.lean}
\lean{LanglandsFirstMainLemma.traceIdealLowerBound_of_integralGenerator}
\lean{LanglandsFirstMainLemma.trace_lattice_surjective_of_integralGenerator}
\lean{LanglandsFirstMainLemma.ramificationRatioUnit}
\lean{LanglandsFirstMainLemma.coe_ramificationRatioUnit}
\lean{LanglandsFirstMainLemma.ord_ramificationRatioUnit}
\lean{LanglandsFirstMainLemma.ramificationRatioUnit_mem}
\lean{LanglandsFirstMainLemma.ramificationUnitGradedPreHom}
\lean{LanglandsFirstMainLemma.ramificationRatioUnit_mem_succ_iff}
\lean{LanglandsFirstMainLemma.ramificationUnitGradedHom}
\lean{LanglandsFirstMainLemma.ramificationUnitGradedHom_mk}
\lean{LanglandsFirstMainLemma.ramificationUnitGradedHom_injective}
\lean{LanglandsFirstMainLemma.norm_ramificationRatioUnit_eq_one}
\lean{LanglandsFirstMainLemma.residueCharacteristic_eq_degree_of_positive_isLowerBreak}
\lean{LanglandsFirstMainLemma.wild_norm_one_add_sub_one_mem_lattice}
\lean{LanglandsFirstMainLemma.wild_norm_one_add_sub_one_sub_norm_mem_lattice}
\lean{LanglandsFirstMainLemma.wild_norm_critical_sub_trace_sub_norm_mem_lattice}
\lean{LanglandsFirstMainLemma.wild_norm_one_add_sub_one_sub_trace_mem_lattice}
\lean{LanglandsFirstMainLemma.NormMapsUnitFiltration}
\lean{LanglandsFirstMainLemma.normMapsUnitFiltration_zero}
\lean{LanglandsFirstMainLemma.normBelowBreakUnitFiltrationHom}
\lean{LanglandsFirstMainLemma.coe_normBelowBreakUnitFiltrationHom}
\lean{LanglandsFirstMainLemma.normUnitFiltrationQuotient}
\lean{LanglandsFirstMainLemma.normUnitFiltrationQuotient_mk}
\lean{LanglandsFirstMainLemma.normUnitGraded}
\lean{LanglandsFirstMainLemma.normUnitGraded_mk}
\lean{LanglandsFirstMainLemma.exists_normUnitFiltration_oneStepCorrection}
\lean{LanglandsFirstMainLemma.mem_unitFiltration_succ_of_norm_of_graded_injective}
\lean{LanglandsFirstMainLemma.exists_normUnitFiltration_finiteCorrection}
\lean{LanglandsFirstMainLemma.mem_unitFiltration_of_norm_mem_of_graded_injective}
\lean{LanglandsFirstMainLemma.normUnitFiltrationInterval}
\lean{LanglandsFirstMainLemma.normUnitFiltrationInterval_mk}
\lean{LanglandsFirstMainLemma.normUnitFiltrationInterval_surjective}
\lean{LanglandsFirstMainLemma.normUnitFiltrationInterval_injective}
\lean{LanglandsFirstMainLemma.normUnitFiltrationIntervalEquiv}
\lean{LanglandsFirstMainLemma.normUnitFiltrationIntervalEquiv_apply}
\lean{LanglandsFirstMainLemma.residueCard_eq_of_residueDegree_eq_one}
\lean{LanglandsFirstMainLemma.unitGradedPiece_card_eq_of_residueDegree_eq_one}
\lean{LanglandsFirstMainLemma.unitGradedMap_bijective_of_injective_of_residueDegree_eq_one}
\lean{LanglandsFirstMainLemma.wild_normMapsUnitFiltration}
\lean{LanglandsFirstMainLemma.normUnitGraded_injective_belowBreak}
\lean{LanglandsFirstMainLemma.normUnitGraded_bijective_belowBreak}
\lean{LanglandsFirstMainLemma.normUnitGradedEquiv_belowBreak}
\lean{LanglandsFirstMainLemma.exists_wild_norm_aboveBreak_oneStep}
\lean{LanglandsFirstMainLemma.wild_norm_unitFiltration_surjective_break_succ}
\lean{LanglandsFirstMainLemma.ramificationUnitGradedHom_le_normUnitGraded_ker}
\lean{LanglandsFirstMainLemma.exists_exactNormOneRepresentative_of_criticalKernel}
\lean{LanglandsFirstMainLemma.exists_hilbert90_coboundary_of_criticalKernel}
\lean{LanglandsFirstMainLemma.gradedNorm_kernel_positiveBreak}
\lean{LanglandsFirstMainLemma.norm_one_add_sub_one_mem_lattice_one_of_residueDegree_eq_one}
\lean{LanglandsFirstMainLemma.tame_norm_one_add_sub_one_sub_trace_mem_lattice}
\lean{LanglandsFirstMainLemma.normMapsUnitFiltration_one_of_residueDegree_eq_one}
\lean{LanglandsFirstMainLemma.exists_tame_norm_aboveBreak_oneStep}
\lean{LanglandsFirstMainLemma.tame_norm_unitFiltration_surjective_break_succ}
\lean{LanglandsFirstMainLemma.tame_ramificationUnitGradedHom_le_normUnitGraded_ker}
\lean{LanglandsFirstMainLemma.exists_tame_exactNormOneRepresentative_of_criticalKernel}
\lean{LanglandsFirstMainLemma.exists_tame_hilbert90_coboundary_of_criticalKernel}
\lean{LanglandsFirstMainLemma.gradedNorm_kernel_tameBreak}
\lean{LanglandsFirstMainLemma.normMapsUnitFiltration_atBreak}
\lean{LanglandsFirstMainLemma.normMapsUnitFiltration_break_succ}
\lean{LanglandsFirstMainLemma.criticalGradedNorm}
\lean{LanglandsFirstMainLemma.criticalGradedNorm_mk}
\lean{LanglandsFirstMainLemma.criticalGradedNorm_kernel_card_positiveBreak}
\lean{LanglandsFirstMainLemma.ramification_gradedNorm_mulExact_positiveBreak}
\lean{LanglandsFirstMainLemma.criticalGradedNorm_cokernel_card_positiveBreak}
\lean{LanglandsFirstMainLemma.criticalGradedNorm_range_card_mul_degree}
\lean{LanglandsFirstMainLemma.gradedNorm_kernel}
\lean{LanglandsFirstMainLemma.ramification_criticalGradedNorm_mulExact}
\lean{LanglandsFirstMainLemma.criticalGradedNorm_kernel_card}
\lean{LanglandsFirstMainLemma.criticalGradedNorm_cokernel_card}
\lean{LanglandsFirstMainLemma.criticalNormLinearCoefficient}
\lean{LanglandsFirstMainLemma.criticalNormRamificationLambda}
\lean{LanglandsFirstMainLemma.criticalNormPolynomialValue}
\lean{LanglandsFirstMainLemma.criticalNormPolynomialValue_exact_of_positiveBreak}
\lean{LanglandsFirstMainLemma.criticalNormLinearCoefficient_eq_neg_ramificationLambda_pow}
\lean{LanglandsFirstMainLemma.criticalNormGaloisZModEquiv}
\lean{LanglandsFirstMainLemma.criticalNormGaloisZModEquiv_apply_val}
\lean{LanglandsFirstMainLemma.criticalNormResidueDisplacement_zmod}
\lean{LanglandsFirstMainLemma.normMapsUnitFiltration_belowBreak}
\lean{LanglandsFirstMainLemma.gradedNorm_bijective_belowBreak}
\lean{LanglandsFirstMainLemma.gradedNorm_kernel_belowBreak}
\lean{LanglandsFirstMainLemma.gradedNorm_range_belowBreak}
\lean{LanglandsFirstMainLemma.ramification_gradedNorm_mulExact_belowBreak}
\lean{LanglandsFirstMainLemma.gradedNormEquiv_belowBreak}
\lean{LanglandsFirstMainLemma.normUnitFiltrationEquiv_belowBreak}
\lean{LanglandsFirstMainLemma.exists_norm_belowBreak_finiteCorrection}
\lean{LanglandsFirstMainLemma.mem_unitFiltration_of_norm_mem_belowBreak}
\lean{LanglandsFirstMainLemma.exists_norm_representative_mod_break}
\lean{LanglandsFirstMainLemma.mem_unitFiltration_break_of_norm_mem}
\lean{LanglandsFirstMainLemma.normFieldFiltrationQuotient}
\lean{LanglandsFirstMainLemma.normFieldFiltrationQuotient_mk}
\lean{LanglandsFirstMainLemma.normFieldFiltrationEquiv_belowBreak}
\leanok
\SourceText{Theorems~\ref{thm:graded-norm} and
\ref{thm:norm-filtration}, including the below-break and critical-break
parts of their proofs.}
\uses{mod:ramification-lowergroups,mod:ramification-herbrand,mod:localfield-successivelifting}
\FormalizationContract The exact quotient norm keeps its numerator and
denominator containments separate.  The integral uniformizer, its
$\mathcal O_F$-algebra generator proof, the residue-degree-one hypothesis,
and the unique-break witness remain explicit inputs.  No declaration imports
or assumes the final norm-filtration theorem or a graded-kernel theorem.
The tame coordinate and positive-depth coordinate are not conflated.
Approximate critical norm-one classes are corrected to exact norm one before
Hilbert 90, whose quotient orientation is $y/g(y)$.  Finite below-break
splicing uses only $i<t$, critical quotients use the denominator $t+1$, and
the first above-break correction uses the distinct Herbrand depth
$\Psi(t+1)=t+\ell$.  Kernel and image directions, representative
independence, and the finite kernel, range, and cokernel cardinalities are
all stated as quotient-group facts.  The critical polynomial and its
coefficient identity are stated for an arbitrary positive break $t$, with
the residue-field transport and the integral-generator witness explicit.
\end{theorem}
